A seguir, descreve-se a metodologia utilizada para a construção, ajuste e avaliação do sistema fuzzy Takagi-Sugeno aplicado à aproximação da função alvo $f(x) = e^{-x/5} \cdot \sin(3x) + 0.5 \cdot \sin(x)$ no intervalo $x \in [0, 10]$.

\subsection{Geração do Conjunto de Dados}

Foram gerados 200 pontos uniformemente espaçados no intervalo $[0, 10]$. Para cada ponto $x_k$, foi calculado o valor da função alvo $f(x_k)$. Esses dados foram utilizados tanto para o treinamento quanto para a avaliação dos modelos fuzzy.

\subsection{Estrutura do Sistema Fuzzy Takagi-Sugeno}

O sistema fuzzy foi estruturado com uma variável de entrada $x$ e uma saída $y \approx f(x)$. Foram consideradas três categorias para $x$ (Baixo, Médio, Alto), implementadas por meio de funções de pertinência gaussianas, triangulares e trapezoidais. O número de funções de pertinência (regras) foi variado entre 5 e 20, permitindo analisar o impacto da granularidade do particionamento do universo de entrada na qualidade da aproximação.

\subsection{Takagi-Sugeno de Ordem Zero}

No modelo de ordem zero, cada regra fuzzy possui um consequente constante $c_i$. A saída do sistema é calculada como a média ponderada dos consequentes pelas pertinências:

\[
y(x) = \frac{\sum_{i=1}^{n} \mu_i(x) \cdot c_i}{\sum_{i=1}^{n} \mu_i(x)}
\]

Os consequentes $c_i$ são inicialmente determinados pela média ponderada dos valores da função alvo, considerando a pertinência de cada regra.

\subsection{Takagi-Sugeno de Primeira Ordem}

No modelo de primeira ordem, cada regra fuzzy possui um consequente linear $y = a_i x + b_i$. Os parâmetros $a_i$ e $b_i$ são ajustados por mínimos quadrados ponderados, utilizando os graus de pertinência como pesos. A saída global é dada por:

\[
y(x) = \frac{\sum_{i=1}^{n} \mu_i(x) \cdot (a_i x + b_i)}{\sum_{i=1}^{n} \mu_i(x)}
\]

\subsection{Avaliação de Desempenho}

O desempenho dos modelos foi avaliado utilizando o erro quadrático médio (RMSE) entre a saída do sistema fuzzy e a função alvo. O RMSE foi calculado para cada configuração de função de pertinência, para ambos os modelos (ordem zero e primeira ordem), e para ambos os cenários de número de regras ($n=5$ e $n=20$).

\subsection{Visualização dos Resultados}

Foram gerados gráficos comparando a função alvo com as aproximações obtidas pelos modelos Takagi-Sugeno de ordem zero e primeira ordem, para cada tipo de função de pertinência e para cada quantidade de regras. Também foram plotados os erros antes e depois da otimização dos consequentes.

\subsection{Otimização dos Consequentes (Gradiente Descendente)}

Os consequentes do modelo de ordem zero foram otimizados por gradiente descendente, minimizando o RMSE. A atualização dos consequentes foi feita iterativamente, ajustando cada $c_i$ de acordo com o gradiente do erro ponderado pelas pertinências.

\subsection{Análise Comparativa}

Os resultados evidenciam a eficiência do sistema fuzzy Takagi-Sugeno na aproximação de funções não lineares, ressaltando a influência do número de funções de pertinência (regras), dos parâmetros do modelo e das técnicas de ajuste dos consequentes na obtenção de aproximações mais precisas. A comparação entre $n=5$ e $n=20$ regras permitiu observar que o aumento do número de funções de pertinência geralmente resulta em melhor desempenho (menor RMSE), ao custo de maior complexidade do modelo.